\documentclass[11pt]{article}

\usepackage[margin=1in]{geometry}
\usepackage{amsmath,amssymb,amsthm}
\usepackage{booktabs}
\usepackage{hyperref}
\usepackage{xcolor}
\usepackage{listings}

\hypersetup{
  colorlinks=true,
  linkcolor=blue,
  citecolor=blue,
  urlcolor=blue
}

\lstdefinelanguage{Coq}{
  keywords={Definition,Theorem,Lemma,Proof,Qed,Record,Fixpoint,forall,fun,match,with,end,Prop,if,then,else,Require,Import,Open,Scope},
  keywordstyle=\color{blue},
  comment=[s]{(*}{*)},
  commentstyle=\color{gray},
  morecomment=[s]{(**}{*)},
  basicstyle=\ttfamily\small,
  columns=fullflexible,
  breaklines=true
}

\lstdefinelanguage{bash}{
  keywords={cd,make,python3,pdflatex,bibtex},
  keywordstyle=\color{blue},
  basicstyle=\ttfamily\small,
  columns=fullflexible,
  breaklines=true
}

\newtheorem{theorem}{Theorem}
\newtheorem{definition}{Definition}
\newtheorem{lemma}{Lemma}

\title{Machine-Checked Verification of Crypto Option Hedged Portfolios under Discontinuous Price Shocks}
\author{FinShockVerify Project}
\date{\today}

\begin{document}

\maketitle

\begin{abstract}
Crypto portfolios hedged with options are exposed to abrupt spot-price drops,
volatile derivatives markets, venue risk, custody risk, and settlement
ambiguity. A protective put can provide a mathematical downside floor, but only
under explicit assumptions about payoff definition, enforceability, settlement,
quantity matching, and premium accounting. This paper presents FinShockVerify,
a prototype Rocq/Coq-based framework for machine-checking the accounting and
risk invariants of crypto option hedged portfolios under discontinuous price
shocks. The current artifact formalizes a single-asset portfolio state,
price-jump transitions, fee-aware trades, finite weighted stress scenarios, a
loss-floor checker, and a crypto protective-put payoff model using exact
rational arithmetic. It proves that trades are self-financing up to explicit
fees, that price jumps produce the expected mark-to-market profit and loss,
that the executable loss-floor checker is sound with respect to the formal
inequality, that a long-only fully invested portfolio cannot lose more than a
common asset-return floor, and that a matching long-crypto plus long-put
position has terminal wealth bounded below by cash minus premium plus units
times strike. The contribution is not a market-prediction model or legal
opinion. It is an artifact and design pattern for audit-grade verification of
explicitly specified crypto portfolio stress and option-hedge logic.
\end{abstract}

\noindent\textbf{Keywords:} formal verification, Rocq, Coq, crypto options,
protective puts, portfolio risk, stress testing, price jumps, model validation,
financial technology

\section{Introduction}

Financial portfolios are exposed not only to gradual price changes but also to
discontinuous jumps: earnings surprises, macroeconomic announcements, liquidity
events, defaults, devaluations, market halts, and sudden repricings of risk.
Portfolio risk systems therefore encode rules for marking positions to market,
applying stress scenarios, rebalancing positions, charging transaction costs,
and checking constraints such as solvency, leverage, mandate limits, or loss
floors. These rules are often implemented in spreadsheets, scripts, risk
engines, or portfolio-management platforms. The mathematical formulas may be
simple, but the surrounding state transitions are a common source of errors:
sign mistakes, inconsistent treatment of cash, incorrect fee accounting,
incorrect pass/fail thresholds, and mismatch between the risk policy and its
implementation.

This paper studies a narrow question:

\begin{quote}
Can a proof assistant be used to verify that crypto option hedge logic
preserves specified downside floors, accounting invariants, and stress-test
claims under discontinuous price shocks?
\end{quote}

We answer this question with FinShockVerify, a small open-source prototype
implemented in Rocq, the proof assistant historically known as Coq. The
prototype intentionally begins with exact rational arithmetic and finite stress
scenarios. This is a design choice. Continuous-time jump processes, stochastic
calculus, and floating-point verification are important, but they are not the
most direct path to a useful first verifier. Many operational risk checks are
already defined over finite shock tables: equity down 20 percent, credit spreads
up, rates down, foreign exchange up, or sector-specific shocks. Such scenarios
are discrete, auditable, and close to the workflows used by risk teams.

The main contributions of this paper are:

\begin{enumerate}
  \item A formal model of portfolio state for cash, risky holdings, prices,
  price jumps, and fee-aware trades.
  \item A formal payoff model for a crypto protective put and a proof that a
  matching long-spot plus long-put position has a terminal value floor at the
  option strike, net of explicitly modeled premium and cash.
  \item Machine-checked proofs of basic accounting invariants, including
  self-financing trades with explicit fees and mark-to-market profit and loss
  under a price jump.
  \item A finite weighted stress-scenario model with exact rational arithmetic.
  \item A machine-checked soundness theorem for a Boolean loss-floor checker.
  \item A proof that long-only fully invested portfolios inherit a common
  asset-return floor at the portfolio level.
  \item A small command-line checker that mirrors the formal definitions over
  CSV inputs and demonstrates a path toward practical model-validation
  tooling.
\end{enumerate}

The contribution is deliberately modest but concrete. FinShockVerify does not
claim to prove that a portfolio is safe, profitable, or correctly priced. It
claims something more precise: given an explicitly specified model of portfolio
state and stress rules, the proof assistant can certify that the executable
checks and transitions satisfy the stated mathematical invariants.

\section{Motivation}

Formal verification is already used in domains where small implementation
errors can create large losses: compilers, cryptographic protocols, operating
systems, smart contracts, and safety-critical systems. Financial risk
infrastructure shares several properties with those domains. It uses precise
rules, operates over state, produces decisions that affect money and capital,
and must be explainable to auditors, regulators, clients, or counterparties.

In financial institutions, model risk management guidance emphasizes robust
model development, validation, governance, and controls. The Federal Reserve's
SR 11-7 guidance describes model risk as the risk of adverse consequences from
decisions based on incorrect or misused model outputs, and lists uses including
stress testing, risk measurement, compliance with internal limits, and financial
reporting \cite{fed2011sr117}. The Basel market-risk framework likewise places
stress and model assessment at the center of market-risk capital, including
expected shortfall models and separate treatment for non-modellable risk
factors \cite{basel2019marketrisk}. These frameworks do not require proof
assistants, but they create a natural environment for proof-carrying
validation artifacts: if an institution says that a risk rule has a precise
meaning, a machine-checked proof can help establish that the implementation
matches that meaning.

The motivating example for this prototype is a long-only equity-bond portfolio:

\begin{center}
\begin{tabular}{lrr}
\toprule
Asset & Weight & Stress return \\
\midrule
Equity & 70\% & -20\% \\
Bonds & 30\% & 3\% \\
\bottomrule
\end{tabular}
\end{center}

The portfolio stress return is
\[
  0.70 \cdot (-0.20) + 0.30 \cdot 0.03 = -0.131,
\]
or -13.1 percent. If the loss floor is -15 percent, the scenario passes. A
unit test can check this example. A proof assistant can do more: it can prove
that the checker used for this computation is sound for every finite scenario
in the model, and that entire classes of portfolios satisfy general bounds.

\section{System Overview}

FinShockVerify is organized around a small formal core and an executable
adapter. The formal core is written in Rocq. The current repository contains
four proof modules:

\begin{itemize}
  \item \texttt{theories/Core/SingleAsset.v}: portfolio value, price jumps,
  and fee-aware trades for one risky asset.
  \item \texttt{theories/Risk/WeightedStress.v}: finite weighted stress
  scenarios, net weight, gross leverage, stress return, and loss-floor checks.
  \item \texttt{theories/Options/CryptoProtectivePut.v}: crypto spot plus
  long-put payoff logic and protective-put floor theorems.
  \item \texttt{theories/Examples/LongOnlyETF.v}: a checked equity-bond
  example.
\end{itemize}

The executable prototype is \texttt{tools/stress\_check.py}, a Python script
that reads CSV scenario files and computes net weight, gross leverage, stress
return, and pass/fail status using exact fractions. The Python tool is not
itself formally verified in the current artifact. It is included to show the
intended product direction: verified specifications should be usable through
normal operational interfaces such as CSV files, scenario tables, and audit
reports.

\section{Formal Model}

The MVP begins with a one-asset state. This is not intended to be a complete
portfolio model. It is the smallest useful state transition on which price
jump and trade accounting can be made precise.

\subsection{State and Value}

\begin{definition}[Portfolio state]
A state consists of cash, risky-asset shares, and the current price:
\[
  s = (\mathrm{cash}, \mathrm{shares}, \mathrm{price}).
\]
\end{definition}

The portfolio value is
\[
  V(s) = \mathrm{cash}(s) + \mathrm{shares}(s)\cdot\mathrm{price}(s).
\]

The Rocq definition is:

\begin{lstlisting}[language=Coq]
Record State := {
  cash : Q;
  shares : Q;
  price : Q
}.

Definition portfolio_value (s : State) : Q :=
  cash s + shares s * price s.
\end{lstlisting}

All quantities use Rocq rational numbers \texttt{Q}. This gives exact
arithmetic and avoids the ambiguity of floating-point rounding. Later versions
can add verified floating-point bounds, but the first layer should be a clear
mathematical specification.

\subsection{Price Jumps}

A discontinuous price jump is modeled as a return \(r\) applied to the current
price:
\[
  P' = P(1+r).
\]
Cash and shares do not change during the mark-to-market jump:

\begin{lstlisting}[language=Coq]
Definition apply_price_return (r : Q) (s : State) : State :=
  {|
    cash := cash s;
    shares := shares s;
    price := price s * (1 + r)
  |}.
\end{lstlisting}

The first price-jump theorem states that the change in portfolio value equals
the position size times the price change:

\begin{theorem}[Price-jump profit and loss]
For every state \(s\) and return \(r\),
\[
  V(\mathrm{jump}(r,s)) - V(s)
  = \mathrm{shares}(s)\cdot\mathrm{price}(s)\cdot r.
\]
\end{theorem}

In Rocq:

\begin{lstlisting}[language=Coq]
Theorem price_jump_pnl :
  forall s r,
    portfolio_value (apply_price_return r s) - portfolio_value s
    == shares s * price s * r.
\end{lstlisting}

\subsection{Trades and Fees}

A trade changes shares and cash at the current price. Buying \(q\) shares costs
\(qP\); selling corresponds to negative \(q\). A nonnegative fee \(f\) is
deducted from cash:

\[
  \mathrm{cash}' = \mathrm{cash} - qP - f,
  \qquad
  \mathrm{shares}' = \mathrm{shares} + q.
\]

\begin{lstlisting}[language=Coq]
Definition trade (quantity fee : Q) (s : State) : State :=
  {|
    cash := cash s - quantity * price s - fee;
    shares := shares s + quantity;
    price := price s
  |}.
\end{lstlisting}

The basic accounting theorem is:

\begin{theorem}[Self-financing trade with explicit fee]
For every state \(s\), quantity \(q\), and fee \(f\),
\[
  V(\mathrm{trade}(q,f,s)) = V(s) - f.
\]
\end{theorem}

In Rocq:

\begin{lstlisting}[language=Coq]
Theorem trade_self_financing :
  forall s quantity fee,
    portfolio_value (trade quantity fee s) == portfolio_value s - fee.
\end{lstlisting}

This theorem is simple, but it captures a critical modeling principle: trading
does not create or destroy portfolio value except through explicitly modeled
costs. In larger systems, this is the type of invariant that can prevent
subtle cash-position errors.

\section{Crypto Protective-Put Hedge}

Crypto option hedges are attractive because large spot-price moves are common
and because listed or centrally intermediated derivatives can offer a more
standardized risk-management surface than ad hoc off-chain agreements. They
also introduce assumptions that must be made explicit. The CFTC states that
virtual currencies such as Bitcoin have been treated as commodities under the
Commodity Exchange Act and that its jurisdiction is implicated when virtual
currency is used in a derivatives contract \cite{cftcBitcoinBasics}. Its
customer advisory notes that futures and options may be used by hedgers to
protect against virtual-currency volatility, while warning about volatility,
cash-market manipulation, cyber risk, margin amplification, and the importance
of understanding cash-settlement indexes \cite{cftcVirtualCurrencyAdvisory}.
The CFTC's Coinflip action also illustrates that platforms offering Bitcoin
options can raise registration and market-facility issues
\cite{cftcCoinflip2015}.

For this reason, FinShockVerify separates the mathematical hedge theorem from
external regulatory and operational assumptions. The proof below assumes that
the put option is valid, enforceable, quantity-matched, and settles according
to its payoff definition. It does not prove venue registration, custody safety,
oracle integrity, counterparty solvency, or legal eligibility.

\subsection{Put Payoff}

For final spot price \(S_T\) and strike \(K\), a long put has payoff
\[
  (K - S_T)^+ = \max(K - S_T, 0).
\]

The Rocq model uses an exact rational positive-part function:

\begin{lstlisting}[language=Coq]
Definition positive_part (x : Q) : Q :=
  match Qlt_le_dec 0 x with
  | left _ => x
  | right _ => 0
  end.

Definition long_put_payoff (final_spot strike : Q) : Q :=
  positive_part (strike - final_spot).
\end{lstlisting}

\subsection{Protected Terminal Wealth}

For \(q\) units of crypto, cash \(C\), option premium \(P\), final spot
\(S_T\), and strike \(K\), the protected terminal wealth is modeled as:
\[
  W = C - P + q\left(S_T + (K - S_T)^+\right).
\]

In Rocq:

\begin{lstlisting}[language=Coq]
Definition protected_crypto_value
    (units final_spot strike : Q) : Q :=
  units * (final_spot + long_put_payoff final_spot strike).

Definition protected_crypto_wealth
    (cash premium units final_spot strike : Q) : Q :=
  cash - premium + protected_crypto_value units final_spot strike.
\end{lstlisting}

The central hedge theorem is:

\begin{theorem}[Protective-put terminal wealth floor]
If \(q \geq 0\), then for every final spot price \(S_T\),
\[
  C - P + qK
  \leq
  C - P + q\left(S_T + (K - S_T)^+\right).
\]
\end{theorem}

In Rocq:

\begin{lstlisting}[language=Coq]
Theorem protected_crypto_wealth_floor :
  forall cash premium units final_spot strike,
    0 <= units ->
    cash - premium + units * strike
    <= protected_crypto_wealth cash premium units final_spot strike.
\end{lstlisting}

This is the precise version of the ``harsh drop'' safety claim. If the crypto
price falls from \(S_0\) to zero, or to any lower stressed value, the terminal
spot-plus-put component remains at least \(qK\), assuming the put settles as
specified. The proof does not remove the option premium; it includes it
explicitly. Therefore the guaranteed net floor is \(C - P + qK\), not \(qK\).

The artifact also proves that the protected position dominates unhedged spot
exposure:

\begin{lstlisting}[language=Coq]
Theorem protected_crypto_dominates_spot :
  forall units final_spot strike,
    0 <= units ->
    crypto_spot_value units final_spot
    <= protected_crypto_value units final_spot strike.
\end{lstlisting}

Finally, when the final spot price is below the strike, the protected value is
exactly the strike floor:

\begin{lstlisting}[language=Coq]
Theorem protected_crypto_value_equals_floor_below_strike :
  forall units final_spot strike,
    final_spot <= strike ->
    protected_crypto_value units final_spot strike == units * strike.
\end{lstlisting}

\subsection{Safety Case Boundaries}

The proof should be used as part of a safety case, not as a standalone
guarantee. A complete crypto option hedge validation should include at least
the following assumptions:

\begin{itemize}
  \item the option contract is legally valid for the relevant participant and
  jurisdiction;
  \item the venue or intermediary is appropriately registered or exempt;
  \item the settlement index and final settlement price are defined and
  auditable;
  \item the put quantity matches the spot crypto exposure;
  \item option premium, fees, margin, and funding are included in wealth;
  \item custody, wallet, exchange, counterparty, and cyber risks are modeled
  separately;
  \item stress scenarios include gap moves, flash crashes, liquidity shocks,
  and settlement/index disruption.
\end{itemize}

\section{Finite Stress Scenarios}

The second part of the MVP formalizes finite stress scenarios. Each leg has a
portfolio weight \(w_i\) and an asset return \(r_i\). The stress return is the
weighted sum:
\[
  R = \sum_i w_i r_i.
\]

The Rocq record is:

\begin{lstlisting}[language=Coq]
Record WeightedShock := {
  weight : Q;
  asset_return : Q
}.
\end{lstlisting}

\subsection{Net Weight, Gross Leverage, and Stress Return}

FinShockVerify defines:

\[
  \mathrm{netWeight} = \sum_i w_i,
  \qquad
  \mathrm{grossLeverage} = \sum_i |w_i|,
  \qquad
  \mathrm{stressReturn} = \sum_i w_i r_i.
\]

In Rocq, these are recursive functions over a list of \texttt{WeightedShock}
records:

\begin{lstlisting}[language=Coq]
Fixpoint net_weight (scenario : list WeightedShock) : Q :=
  match scenario with
  | [] => 0
  | leg :: rest => weight leg + net_weight rest
  end.

Fixpoint gross_leverage (scenario : list WeightedShock) : Q :=
  match scenario with
  | [] => 0
  | leg :: rest => Qabs (weight leg) + gross_leverage rest
  end.

Fixpoint stress_return (scenario : list WeightedShock) : Q :=
  match scenario with
  | [] => 0
  | leg :: rest => weight leg * asset_return leg + stress_return rest
  end.
\end{lstlisting}

The model also defines common portfolio conditions:

\begin{lstlisting}[language=Coq]
Definition long_only (scenario : list WeightedShock) : Prop :=
  Forall (fun leg => 0 <= weight leg) scenario.

Definition returns_floor (floor : Q) (scenario : list WeightedShock) : Prop :=
  Forall (fun leg => floor <= asset_return leg) scenario.

Definition fully_invested (scenario : list WeightedShock) : Prop :=
  net_weight scenario == 1.
\end{lstlisting}

\subsection{Executable Loss-Floor Check}

A practical stress checker often needs to answer a Boolean question: did this
scenario pass a loss floor?

\[
  \mathrm{passesLossFloor}(F, S) =
  \begin{cases}
  \mathrm{true} & \text{if } F \leq \mathrm{stressReturn}(S),\\
  \mathrm{false} & \text{otherwise.}
  \end{cases}
\]

The Rocq definition is:

\begin{lstlisting}[language=Coq]
Definition passes_loss_floor (floor : Q) (scenario : list WeightedShock) : bool :=
  Qle_bool floor (stress_return scenario).
\end{lstlisting}

The soundness theorem connects the Boolean checker to the mathematical
inequality:

\begin{theorem}[Loss-floor checker soundness]
If \texttt{passes\_loss\_floor floor scenario = true}, then
\[
  \mathrm{floor} \leq \mathrm{stressReturn}(\mathrm{scenario}).
\]
\end{theorem}

In Rocq:

\begin{lstlisting}[language=Coq]
Theorem passes_loss_floor_sound :
  forall floor scenario,
    passes_loss_floor floor scenario = true ->
    floor <= stress_return scenario.
\end{lstlisting}

This theorem matters because it ties executable pass/fail behavior to a formal
specification. It is a small version of a larger product goal: risk reports
should be accompanied by proof artifacts that explain what their pass/fail
status means.

\subsection{Long-Only Loss Floor}

The main general theorem in the current MVP is a long-only loss-floor result.
Suppose every portfolio weight is nonnegative, the weights sum to one, and
every asset return is at least \(-\epsilon\). Then the portfolio return is also
at least \(-\epsilon\).

\begin{theorem}[Long-only fully invested loss floor]
For every finite scenario \(S\) and rational \(\epsilon \geq 0\), if
\[
  w_i \geq 0,\quad \sum_i w_i = 1,\quad r_i \geq -\epsilon
\]
for every leg \(i\), then
\[
  \sum_i w_i r_i \geq -\epsilon.
\]
\end{theorem}

In Rocq:

\begin{lstlisting}[language=Coq]
Theorem long_only_fully_invested_loss_floor :
  forall scenario eps,
    0 <= eps ->
    long_only scenario ->
    returns_floor (- eps) scenario ->
    fully_invested scenario ->
    - eps <= stress_return scenario.

Theorem protected_crypto_wealth_floor :
  forall cash premium units final_spot strike,
    0 <= units ->
    cash - premium + units * strike
    <= protected_crypto_wealth cash premium units final_spot strike.
\end{lstlisting}

The theorem is mathematically elementary. Its significance is architectural:
it demonstrates how portfolio mandates and stress assumptions can be stated as
machine-checkable hypotheses and connected to portfolio-level risk guarantees.
Later versions can generalize this pattern to leverage caps, position limits,
sector constraints, margin rules, liquidity haircuts, and expected shortfall
over finite scenario distributions.

\section{Prototype Implementation}

\subsection{Rocq Artifact}

The Rocq artifact compiles using a standard \texttt{\_CoqProject} file and a
Makefile. The build command is:

\begin{lstlisting}[language=bash]
make build
\end{lstlisting}

The current checked source files are summarized in Table~\ref{tab:artifact}.

\begin{table}[h]
\centering
\begin{tabular}{lll}
\toprule
File & Role & Selected checked statements \\
\midrule
\texttt{SingleAsset.v} & Core state transitions &
\texttt{trade\_self\_financing}, \texttt{price\_jump\_pnl} \\
\texttt{Vanilla.v} & Vanilla option primitives &
\texttt{option\_payoff\_nonnegative} \\
\texttt{PortfolioValidation.v} & Option validation &
\texttt{covered\_protective\_put\_wealth\_floor} \\
\texttt{WeightedStress.v} & Finite stress scenarios &
\texttt{passes\_loss\_floor\_sound},
\texttt{long\_only\_fully\_invested\_loss\_floor} \\
\texttt{CryptoProtectivePut.v} & Crypto option hedge &
\texttt{protected\_crypto\_wealth\_floor} \\
\texttt{LongOnlyETF.v} & Example portfolio &
\texttt{equity\_bond\_verified\_loss\_floor} \\
\bottomrule
\end{tabular}
\caption{Current FinShockVerify proof artifact.}
\label{tab:artifact}
\end{table}

\subsection{Command-Line Checker}

The Python command-line checker reads a CSV file with columns:
\texttt{asset}, \texttt{weight}, and \texttt{return}. It computes the same
finite-scenario quantities as the Rocq model, using Python's
\texttt{Fraction} type to preserve exact arithmetic.

Example input:

\begin{lstlisting}
asset,weight,return
Equity,70/100,-20/100
Bonds,30/100,3/100
\end{lstlisting}

Example command:

\begin{lstlisting}[language=bash]
python3 tools/stress_check.py examples/long_only_stress.csv --loss-floor=-15/100
\end{lstlisting}

Example output:

\begin{lstlisting}
FinShockVerify stress check
scenario:       examples/long_only_stress.csv
net weight:     1
gross leverage: 1
stress return:  -13.1000% (-131/1000)
loss floor:     -15.0000% (-3/20)
result:         PASS
\end{lstlisting}

The Python script is best viewed as a user-facing adapter to the formal
specification. A future version should generate executable code directly from
Rocq extraction, or use proof-producing checkers so the command-line result is
linked more tightly to the proof artifact.

\section{Evaluation}

The current evaluation is intentionally small. The goal is to show that the
approach can verify meaningful accounting and stress invariants, not to
benchmark a production risk engine.

\subsection{Case Study 1: Equity-Bond Shock}

The checked example uses a two-leg portfolio:

\begin{align*}
  w_{\mathrm{equity}} &= 70/100, &
  r_{\mathrm{equity}} &= -20/100,\\
  w_{\mathrm{bond}} &= 30/100, &
  r_{\mathrm{bond}} &= 3/100.
\end{align*}

The formal example proves:

\begin{itemize}
  \item the weights are fully invested;
  \item the weights are long-only;
  \item each return is at least -20 percent;
  \item the exact stress return is -131/1000;
  \item the Boolean loss-floor checker passes a -15 percent floor;
  \item the general long-only loss-floor theorem applies to the scenario.
\end{itemize}

This example illustrates the two layers of the system. The concrete
calculation verifies one scenario. The general theorem verifies a class of
scenarios satisfying stated hypotheses.

\subsection{Case Study 2: One-Bitcoin Protective Put}

The crypto hedge module contains a simple checked crash example: one unit of
crypto, a put strike of 80,000, final stressed spot of 30,000, and option
premium of 5,000. The theorem proves:

\[
  0 - 5000 + 1\cdot 80000
  \leq
  W(0,5000,1,30000,80000).
\]

That is, even after a severe drop to 30,000, the protected terminal wealth is
bounded below by 75,000 under the payoff-settlement assumptions. The same
theorem also covers a drop to zero or any other final spot price, because the
proof quantifies over all final spots.

\subsection{Case Study 3: Trade Accounting}

The one-asset trade model verifies that a trade at the current price is
self-financing up to explicit fees. This is a foundational invariant for later
rebalancing models. Without such a theorem, a rebalancing procedure could
accidentally create or destroy value through inconsistent cash and holdings
updates.

\subsection{Case Study 4: Price-Jump Mark-to-Market}

The price-jump model verifies that holding \(h\) shares at price \(P\) and
applying return \(r\) changes portfolio value by \(hPr\). This theorem is a
building block for multi-asset mark-to-market logic, stress loss attribution,
and scenario explanation.

\subsection{Limitations of the Current Evaluation}

The current artifact has five major limitations:

\begin{enumerate}
  \item The trade model is single-asset.
  \item The stress model is finite and deterministic.
  \item The Python CLI is specification-aligned but not extracted from Rocq.
  \item The protective-put theorem assumes correct option settlement and does
  not yet model venue default, custody loss, liquidation, margin calls, or
  index manipulation.
  \item The artifact does not yet verify floating-point behavior.
\end{enumerate}

These limitations are acceptable for a first artifact, but they define the
minimum roadmap before claiming production-grade validation.

\section{Business and Validation Use Case}

The practical value of FinShockVerify lies in model validation and audit
evidence. A financial institution or fintech company may already have a policy
such as:

\begin{quote}
For each approved finite stress scenario, the post-shock portfolio must satisfy
loss, leverage, cash, and mandate constraints after applying the documented
rebalancing rule and transaction-cost model.
\end{quote}

Traditional testing can sample cases. A proof assistant can verify universal
statements about all states satisfying the policy hypotheses. The likely first
commercial product is therefore not a trading strategy. It is a validation
tool that produces:

\begin{itemize}
  \item a formal specification of portfolio transition rules;
  \item executable checks over portfolio and scenario files;
  \item machine-checked proof artifacts for core invariants;
  \item human-readable audit reports explaining assumptions and guarantees.
\end{itemize}

This positioning is consistent with model-risk governance: the system verifies
that specified rules are implemented correctly. It does not certify that the
rules are economically sufficient or that the stress scenarios are complete.

\section{Related Work}

\subsection{Formalized Mathematical Finance}

Coelho's Lean 4 mathematical finance library formalizes a broad collection of
results across stochastic calculus, derivative pricing, risk, portfolio, and
fixed-income theory \cite{coelho2026leanfinance}. That work is broader and
more mathematical than FinShockVerify. Its contribution is reusable verified
foundations for mathematical finance. FinShockVerify instead focuses on an
operational layer: state transitions, finite stress scenarios, and audit-style
checks for portfolio logic.

\subsection{Financial Contract Semantics}

The Findel line of work formalizes a domain-specific language for financial
derivatives in Coq and proves properties of financial contracts
\cite{arusoaie2019findel,arusoaie2020findelcert}. This work is close in spirit
because it treats financial artifacts as formal objects with machine-checked
semantics. FinShockVerify differs by focusing on portfolio-level stress and
rebalancing transitions rather than derivative-contract semantics.

Earlier functional-programming work on compositional financial contracts also
showed how financial products can be represented as precise executable
denotations \cite{jones2000composing}. FinShockVerify follows the same broad
philosophy: financial logic should be made explicit, compositional, and
checkable.

\subsection{Verified Market Mechanisms}

Sarswat and Singh formalized properties of financial-market trades and double
auction mechanisms in Coq, including fairness, uniformity, and individual
rationality \cite{sarswat2020verifiedtrades}. That line of work verifies
market mechanisms. FinShockVerify verifies portfolio state and risk logic after
market events are represented as price shocks or trades.

\subsection{Financial Standards and Risk Governance}

ACTUS provides machine-readable algorithmic representations of financial
contracts and emphasizes consistent, precise representation of contract logic
\cite{actusfrf}. FinShockVerify is complementary: a future version could import
standardized contract cash-flow descriptions and verify portfolio-level stress
rules over the resulting exposures.

Regulatory and supervisory sources motivate the need for validation artifacts.
SR 11-7 emphasizes robust model development, validation, governance, and
controls \cite{fed2011sr117}. Basel's market-risk framework emphasizes
expected shortfall models, non-modellable risk factors, and stress-related
assessment \cite{basel2019marketrisk}. FinShockVerify should be understood as
technical infrastructure that can support such validation, not as a replacement
for expert economic judgment.

\section{Roadmap}

The next version of FinShockVerify should add the following features.

\subsection{Multi-Asset Portfolio State}

The one-asset model should be generalized to vectors or finite maps of assets:

\[
  V(s) = \mathrm{cash}(s) + \sum_i h_i(s) P_i(s).
\]

The corresponding theorem should prove that applying a vector of returns
produces:

\[
  V(s') - V(s) = \sum_i h_i P_i r_i.
\]

\subsection{Verified Rebalancing}

The most important extension is verified rebalancing. A rebalancing rule maps
current holdings, prices, and target weights to trades. The core theorem should
state that the rebalance preserves value up to explicit transaction costs:

\[
  V(\mathrm{rebalance}(s)) = V(s) - \mathrm{costs}.
\]

Additional theorems should prove that target-weight constraints, position
limits, cash constraints, and leverage limits hold after rebalancing when the
algorithm returns success.

\subsection{Transaction Costs and Liquidity Haircuts}

The MVP includes a scalar fee. Production risk systems need richer cost
models: proportional costs, fixed ticket costs, bid-ask spread, slippage,
market-impact approximations, and liquidity haircuts. Each model can be
specified as a function and connected to accounting theorems.

\subsection{Crypto Option Hedge Extensions}

The protective-put theorem should be extended into a full crypto option hedge
validator. The next proof obligations are:

\begin{itemize}
  \item cash-settled and physically settled option variants;
  \item option premium, exchange fees, funding costs, and collateral
  requirements;
  \item margin-call and liquidation transitions for leveraged crypto exposure;
  \item settlement index definitions and admissible oracle/index inputs;
  \item counterparty, clearing, or venue-default assumptions represented as
  explicit hypotheses;
  \item portfolio-level hedges with imperfect quantity matching, multiple
  strikes, collars, spreads, and rolling maturities.
\end{itemize}

\subsection{Finite Expected Shortfall}

The finite stress-scenario model should be extended with scenario
probabilities. This enables verification of expected loss and finite expected
shortfall calculations. Expected shortfall and coherent risk measures are
standard tools in modern risk management
\cite{artzner1999coherent,rockafellar2000cvar}. A finite-scenario
formalization remains simpler than continuous probability theory while covering
many historical and hypothetical scenario workflows.

\subsection{Extraction and Proof-Carrying Reports}

The Python CLI should eventually be replaced or supplemented by extracted code
from Rocq. The long-term product should produce a report containing:

\begin{itemize}
  \item input files and hashes;
  \item scenario results;
  \item checked theorem names;
  \item assumptions used;
  \item reproducible build commands.
\end{itemize}

\section{Threats to Validity}

The primary threat is specification error. A proof assistant can verify that an
implementation matches a formal specification, but it cannot guarantee that the
specification captures all economically relevant risks. Scenario selection,
liquidity assumptions, and model boundaries still require expert judgment.

The second threat is abstraction error. The current model uses rational
arithmetic and simplified state transitions. This is appropriate for a clean
specification, but production systems often use decimal conventions,
rounding, settlement delays, tax lots, corporate actions, and floating-point
analytics. Future work must either verify these details or clearly separate
the idealized specification from implementation-specific behavior.

The third threat is artifact integration. The current CLI mirrors the formal
definitions, but it is not extracted from Rocq. A production verifier should
reduce this gap by extracting executable code, generating proof certificates,
or producing machine-checkable traces.

\section{Conclusion}

FinShockVerify demonstrates a practical first step toward proof-carrying
crypto portfolio risk infrastructure. The current artifact formalizes price
jumps, fee-aware trades, finite stress scenarios, long-only loss floors, and a
crypto protective-put floor using exact rational arithmetic in Rocq. It proves
basic accounting and risk invariants and connects them to a small executable
stress checker.

The main lesson is not that proof assistants can predict markets. They cannot.
The lesson is that many parts of financial risk infrastructure are rule-based,
stateful, and mathematically precise. Those parts can be specified and
machine-checked. For crypto option hedged portfolios exposed to sudden price
jumps, this creates a path toward auditable validation artifacts: executable
checks whose meaning is backed by formal proofs and whose assumptions are
explicit enough for model-risk review.

\appendix

\section{Artifact Reproduction}

The artifact can be built with:

\begin{lstlisting}[language=bash]
cd finshock-verify
make build
\end{lstlisting}

The example stress checker can be run with:

\begin{lstlisting}[language=bash]
python3 tools/stress_check.py examples/long_only_stress.csv --loss-floor=-15/100
\end{lstlisting}

The expected output contains:

\begin{lstlisting}
stress return:  -13.1000% (-131/1000)
result:         PASS
\end{lstlisting}

\section{Selected Rocq Statements}

\begin{lstlisting}[language=Coq]
Theorem trade_self_financing :
  forall s quantity fee,
    portfolio_value (trade quantity fee s) == portfolio_value s - fee.

Theorem price_jump_pnl :
  forall s r,
    portfolio_value (apply_price_return r s) - portfolio_value s
    == shares s * price s * r.

Theorem passes_loss_floor_sound :
  forall floor scenario,
    passes_loss_floor floor scenario = true ->
    floor <= stress_return scenario.

Theorem long_only_fully_invested_loss_floor :
  forall scenario eps,
    0 <= eps ->
    long_only scenario ->
    returns_floor (- eps) scenario ->
    fully_invested scenario ->
    - eps <= stress_return scenario.
\end{lstlisting}

\bibliographystyle{plain}
\bibliography{references}

\end{document}
